\documentclass[11pt]{article}
\usepackage{rabbithole}

\phasenumber{12}
\phasetitle{GitHub as a platform}
\phasesubtitle{Everything that is not git: issues, projects, pages, and search done properly.}

\hypersetup{pdftitle={Rabbit Holes, Phase 12: GitHub as a platform},
            pdfauthor={Accelerate, Manipal University Jaipur},
            pdfsubject={git and GitHub handbook}}

\begin{document}
\maketitlepage

\section{Where git stops}

Git is a program for recording versions of a folder. It has no concept of a pull
request, an issue, a review, a label, a milestone, or a discussion. None of those are
in git and none of them ever will be.

Everything in this phase is GitHub's, built on top. That is why the workshop needed
two tools: \co{git} to move commits, \co{gh} to touch anything in this document.

\begin{why}
Keeping the boundary clear pays off constantly. When something behaves oddly, the
first question is which system owns it. A rejected push is git. A pull request that
will not merge is GitHub. A conflict is git. A required review is GitHub.

It also tells you what survives if you leave. Your commits are yours and clone in
full. Your issues, reviews and discussions are GitHub's, and moving them means an
export tool.
\end{why}

\section{Issues}

An issue is a thread with a state, a title, labels, an assignee, and a number.

\begin{shell}
gh issue create --title "Parser fails on blank lines" --body "..."
gh issue list --state open --label bug
gh issue list --assignee "@me"
gh issue view 42
gh issue view 42 --comments
gh issue comment 42 --body "Reproduced on Windows too"
gh issue close 42 --reason "not planned"
gh issue reopen 42
gh issue edit 42 --add-label bug --add-assignee shashwat-deep
gh issue transfer 42 accelerate-muj/other-repo
\end{shell}

\subsection{Linking issues to pull requests}

Put a closing keyword in the pull request description:

\begin{terminal}
Closes #42
Fixes #17, closes #18
Resolves accelerate-muj/other-repo#9
\end{terminal}

When the pull request merges, GitHub closes those issues automatically and records
the link both ways. Six months later the issue shows exactly which change resolved
it, which is the kind of thing you only appreciate when you need it.

\begin{note}
The keyword must be in the pull request body or a commit message, not in a comment
added afterwards. \co{Closes \#42} typed as a comment does nothing except look like
it worked.
\end{note}

\subsection{Templates}

Put templates in \co{.github/ISSUE\_TEMPLATE/}. A YAML form gives you structured
fields rather than a blank box people ignore:

\begin{terminal}
name: Bug report
description: Something does not work
labels: ["bug"]
body:
  - type: input
    id: version
    attributes:
      label: Git version
      description: Output of `git --version`
    validations:
      required: true
  - type: textarea
    id: steps
    attributes:
      label: Steps to reproduce
      value: |
        1.
        2.
    validations:
      required: true
\end{terminal}

Add \co{.github/ISSUE\_TEMPLATE/config.yml} to point people elsewhere for things
that are not bugs:

\begin{terminal}
blank_issues_enabled: false
contact_links:
  - name: Question about the workshop
    url: https://github.com/accelerate-muj/git-started/discussions
    about: Ask here instead of opening an issue
\end{terminal}

\subsection{Labels and milestones}

\begin{shell}
gh label create "good first issue" --color 7057ff --description "Good for newcomers"
gh label list
gh issue list --label "good first issue"
\end{shell}

\co{good first issue} and \co{help wanted} are special: GitHub surfaces repositories
using them in its own discovery pages. If you want outside contributors, they are
worth using accurately rather than decoratively.

Milestones group issues toward a target, and show a progress bar. Useful for a
release or an event, less useful as a permanent bucket.

\section{Projects}

Projects are spreadsheets and boards over issues and pull requests, with custom
fields you define: status, priority, effort, iteration, or anything else.

\begin{shell}
gh project list --owner accelerate-muj
gh project item-list 3 --owner accelerate-muj
gh project item-add 3 --owner accelerate-muj --url <issue-url>
\end{shell}

The useful part is built-in automation: an item moves to Done when its pull request
merges, or to In Progress when someone is assigned. Set it up once and the board
stops needing manual maintenance, which is the only reason boards survive.

\begin{note}
Projects are per user or per organisation, not per repository, so one board can span
several repositories. That is the main reason to use them over a simple label
scheme.
\end{note}

\section{Discussions}

Issues are for things with a resolution. Discussions are for things without one:
questions, ideas, showing off what you built.

Categories can be Q\&A, where answers can be marked accepted, which makes
Discussions a decent support forum for a club or a small project.

\begin{shell}
gh api repos/:owner/:repo/discussions --jq '.[].title'
\end{shell}

Enable them under repository settings. For a workshop repository they are worth
having, because otherwise every question becomes an issue and the issue list stops
being a list of work.

\section{Pages}

Free static hosting from a repository.

\begin{shell}
gh api -X POST repos/:owner/:repo/pages -f source[branch]=main -f source[path]=/docs
gh browse --settings
\end{shell}

Three ways to publish: a branch, a \co{/docs} folder, or a GitHub Actions workflow.
The Actions route is what any static site generator uses, and it is covered in Phase
14.

Add a \co{CNAME} file containing your domain to use a custom one. HTTPS is automatic.

\section{Gists}

A gist is a tiny repository for a snippet. It has full git history and can be cloned,
which people forget.

\begin{shell}
gh gist create script.sh --public --desc "Bisect helper"
gh gist create --filename notes.md -
gh gist list
gh gist edit <id>
gh gist clone <id>
\end{shell}

\begin{trap}
A secret gist is unlisted, not private. Anyone with the URL can read it, and the URL
is not protected. Never put credentials in one.
\end{trap}

\section{Search, properly}

GitHub's search syntax is genuinely powerful and almost nobody learns it.

\subsection{Finding code}

\begin{terminal}
language:python "def parse_sonnet"
repo:accelerate-muj/git-started path:*.tex rerere
org:accelerate-muj filename:.gitattributes
"force-with-lease" NOT test language:markdown
/git\s+rebase\s+--onto/
\end{terminal}

The last one is a regular expression, which the current code search supports. Also
useful: \co{symbol:functionName} finds definitions rather than every mention.

\subsection{Finding issues and pull requests}

\begin{terminal}
is:issue is:open label:bug sort:created-desc
is:pr is:merged author:@me
is:pr review-requested:@me is:open
is:pr is:open draft:false -author:@me
is:issue no:assignee label:"good first issue"
is:pr merged:>2026-08-01
is:issue involves:@me updated:>2026-08-20
\end{terminal}

\co{review-requested:@me is:open} is the one to bookmark. It is your actual queue,
and it is not the same as your notifications.

\subsection{Finding repositories}

\begin{terminal}
stars:>1000 language:rust pushed:>2026-01-01
org:accelerate-muj archived:false
topic:workshop topic:git
\end{terminal}

From the terminal:

\begin{shell}
gh search code "force-with-lease" --owner accelerate-muj
gh search issues --label bug --state open --owner accelerate-muj
gh search prs --review-requested @me --state open
gh search repos --topic workshop --sort stars
\end{shell}

\section{Notifications}

The default is to notify on everything you are subscribed to, which quickly becomes
noise that gets ignored entirely.

\begin{shell}
gh api notifications --jq '.[] | "\(.subject.type)  \(.repository.full_name)  \(.subject.title)"'
gh api -X PUT notifications -f read=true
\end{shell}

Worth doing in the web interface once: set watching to \textbf{Participating and
mentions} for repositories you contribute to occasionally, and \textbf{All activity}
only for ones you maintain. \co{@mentions} and review requests still reach you either
way.

\section{Things that are quietly useful}

\begin{description}
  \item[Keyboard shortcuts] Press \co{?} on any page. \co{t} opens the file finder in
    a repository, \co{l} jumps to a line, \co{s} focuses search, and \co{.} opens the
    repository in a web editor.
  \item[Press the full stop key] Genuinely. \co{.} in any repository opens VS Code in
    the browser, editing that repository. \co{github.dev} in the URL does the same.
  \item[Blame from the browser] \co{b} on a file view. Combined with
    \co{.git-blame-ignore-revs} from Phase 6 it is better than most local tools.
  \item[Permalinks] Press \co{y} on a file view and the URL rewrites to a specific
    commit. Without it, a line link rots the moment the file changes.
  \item[Suggested changes] In a review, use a \co{suggestion} code block and the
    author can apply your edit with one click.
  \item[Draft pull requests] \co{gh pr create -{}-draft} marks work as in progress.
    Reviewers are not requested until you mark it ready.
  \item[CODEOWNERS] Automatically requests review from the right people per path.
    Phase 15.
  \item[Dependabot] Opens pull requests for outdated or vulnerable dependencies. Add
    \co{.github/dependabot.yml}.
  \item[Insights] The repository's Insights tab has contributor graphs, traffic, and
    a dependency graph. Traffic is the only place to see clone and view counts.
\end{description}

\section{Reference}

\vspace{0.5em}
\begin{tabularx}{\linewidth}{@{}L{6.6cm} X@{}}
\toprule
\rhtablehead{Goal} & \rhtablehead{Command or syntax} \\
\midrule
Close an issue when a PR merges & \co{Closes \#42} in the PR body \\
My review queue & \co{is:pr review-requested:@me is:open} \\
Unassigned beginner issues & \co{is:issue no:assignee label:"good first issue"} \\
Search code in an org & \co{gh search code "text" -{}-owner <org>} \\
Open the web editor & press \co{.} in any repository \\
Permalink to a line & press \co{y} on the file view \\
Create a draft PR & \co{gh pr create -{}-draft} \\
Structured issue forms & \co{.github/ISSUE\_TEMPLATE/*.yml} \\
Redirect questions elsewhere & \co{ISSUE\_TEMPLATE/config.yml} \\
\bottomrule
\end{tabularx}
\vspace{0.5em}

\checkyourself{
\begin{enumerate}
  \item Which things in this phase exist in git itself? (Careful.)
  \item Why does \co{Closes \#42} in a comment not close the issue?
  \item What is the difference between a secret gist and a private one?
  \item Write the search that finds open pull requests waiting on your review.
  \item When should something be a Discussion rather than an Issue?
\end{enumerate}}

\vspace{1em}
{\color{rhborder}\rule{\linewidth}{0.8pt}}

\textbf{Next:} Phase 13, \emph{The gh CLI in full}, including \co{gh api}, which
reaches every part of GitHub that has no dedicated command.

\end{document}
